% ==============================================================================
% Related Work
% ==============================================================================

\section{Related Work}
\label{sec:related_work}

\paragraph{Dual-process theory and LLMs.}
The application of dual-process theory to artificial systems has generated
considerable recent interest.  \citet{lampinen2025nature} provide a
comprehensive review of the connections between dual-process frameworks and
modern AI, arguing that the distinction is better understood as a continuum of
processing characteristics rather than a categorical divide.
\citet{hagendorff2023human} demonstrated that GPT-3.5 reproduces human-like
performance on the Cognitive Reflection Test (CRT), including susceptibility to
heuristic lures, while GPT-4 partially overcomes these biases---a pattern
reminiscent of individual differences in human ``cognitive reflection.''
\citet{codaforno2025dual} tested whether prompting strategies that differentially
recruit fast versus slow processing produce distinct internal representations,
finding overlapping but partially separable subspaces.
\citet{ziabari2025reasoning} showed that interpolating between a base model and
a reasoning-fine-tuned model yields a monotonic gradient of behavioral change on
dual-process tasks, supporting a continuous rather than binary characterization.
Our work moves from behavioral characterization to mechanistic analysis,
testing whether these behavioral gradients have identifiable internal
correlates.

\paragraph{Probing LLM internal states.}
Linear probing has become a standard tool for studying what information is
encoded in neural network representations
\citep{hewitt2019control, burns2022discovering}.  Several recent studies probe
representations specifically related to reasoning and self-monitoring.
\citet{zhang2025probing} showed that hidden states in DeepSeek-R1 encode
self-verification signals---the model's internal state ``knows'' whether its
answer is correct before producing it.
\citet{afzal2025knowing} demonstrated that models internally prepare for
difficult outputs several tokens in advance, suggesting a form of prospective
processing.
\citet{fartale2025recall} disentangled attention heads specialized for factual
recall from those involved in multi-step reasoning.
Our probing analysis extends this line of work to the S1/S2 distinction
specifically, asking whether the conflict versus no-conflict condition is
linearly decodable from residual stream activations and whether this
decodability varies systematically across layers and models.

\paragraph{SAE-based analysis of reasoning.}
Sparse autoencoders (SAEs) decompose neural network activations into
interpretable features, offering finer-grained access to model computation
than linear probes.
\citet{goodfire2025r1} identified specific SAE features in DeepSeek-R1
associated with backtracking, self-reference, and answer urgency, providing
the first mechanistic decomposition of chain-of-thought reasoning.
\citet{fang2026sae} demonstrated that steering SAE features can directly
modulate reasoning quality, while \citet{yang2026step} analyzed how SAE
feature activations evolve across reasoning steps.
Critically, \citet{ma2026falsification} showed that 45--90\% of claimed
``reasoning features'' in SAEs are spurious---they activate in response to
specific tokens rather than the reasoning process itself.  We adopt their
falsification framework as a mandatory control throughout our SAE analysis.

\paragraph{Cognitive biases as interpretability benchmarks.}
Cognitive bias tasks are particularly well-suited to mechanistic interpretability
because they provide \emph{contrastive} conditions: the same task structure can
be presented with or without a heuristic lure, enabling clean comparison of
internal processing on matched stimuli.
\citet{rao2024cobbler} benchmarked cognitive biases in LLMs-as-evaluators.
Most relevant to our work, \citet{huang2026cogbias} recently combined linear
probes with activation steering on a cognitive bias benchmark, achieving
26--32\% bias reduction via internal intervention.  Our work differs in three
respects: (i)~we study the nature of the S1/S2-like processing distinction
rather than steering models away from biases; (ii)~we employ five
complementary mechanistic methods rather than probes alone; and (iii)~we
systematically compare standard and reasoning-trained models sharing the
same architecture.

\paragraph{Chain-of-thought faithfulness.}
The reliability of CoT reasoning as a window into model computation is a
prerequisite for interpreting thinking traces.
\citet{lanham2023measuring} found that many CoT steps can be removed without
affecting the final answer, suggesting low faithfulness.
\citet{jiang2025aha} showed that only ${\sim}2.3$\% of reasoning steps
causally influence model outputs, and that apparent self-correction markers
(``wait,'' ``actually'') are largely decorative.
These findings motivate our emphasis on \emph{activation-level} rather than
\emph{token-level} analysis: we extract representations from the residual
stream and attention patterns, not from the content of the generated text.
